
\documentclass{article}
\usepackage{colortbl}
\usepackage{makecell}
\usepackage{multirow}
\usepackage{supertabular}

\begin{document}

\newcounter{utterance}

\centering \large Interaction Transcript for game `codenames', experiment `frequency\_low', episode 6 with reka{-}flash{-}3.1{-}t0.0.
\vspace{24pt}

{ \footnotesize  \setcounter{utterance}{1}
\setlength{\tabcolsep}{0pt}
\begin{supertabular}{c@{$\;$}|p{.15\linewidth}@{}p{.15\linewidth}p{.15\linewidth}p{.15\linewidth}p{.15\linewidth}p{.15\linewidth}}
    \# & $\;$A & \multicolumn{4}{c}{Game Master} & $\;\:$B\\
    \hline

    \theutterance \stepcounter{utterance}  
    & & & \multicolumn{2}{p{0.3\linewidth}}{
        \cellcolor[rgb]{0.9,0.9,0.9}{
            \makecell[{{p{\linewidth}}}]{
                \texttt{\tiny{[GM$|$GM]}}
                {'hidden': {'team': ['chilliness', 'defacement', 'liquidness', 'glasshouse', 'eggroll', 'ungodliness', 'bandmaster', 'unbefitting', 'ninepin'], 'innocent': ['slouchy', 'musicologist', 'cattiness', 'duplicator', 'duckweed', 'bagpiper', 'toxicologist'], 'opponent': ['secession', 'gentlemanliness', 'candleholder', 'shedder', 'inexistence', 'plantain', 'hydroplane', 'hotelkeeper'], 'assassin': ['geophysics']}, 'revealed': {'team': {'team': [], 'innocent': [], 'opponent': [], 'assassin': []}, 'opponent': {'team': [], 'innocent': [], 'opponent': [], 'assassin': []}}}
            }
        }
    }
    & & \\ \\

    \theutterance \stepcounter{utterance}  
    & & \multicolumn{4}{p{0.6\linewidth}}{
        \cellcolor[rgb]{0.9,0.9,0.9}{
            \makecell[{{p{\linewidth}}}]{
                \texttt{\tiny{[P1$\langle$GM]}}
                \texttt{Provide one single{-}word clue that relates to some of your team's words from the list below. You can choose to how many and to which words your clue relates to.} \\
\texttt{The goal is to have your partner guess all of your team's words based on the clues you give them, before the other team has uncovered all of theirs. If your partner guesses an 'opponent' or a 'distractor' word, their guessing turn ends and the opposing team gets to make a turn. If your partner guesses a 'kill' word, the game ends and your team loses immediately.} \\
\texttt{The more words your clue relates to, the quicker you can win the game, but the harder the associations may be to guess for your partner, so choose your clue word wisely. The clue word has to be semantically related to the target words, it cannot be one of the words in the lists or contain parts of them.} \\
\\ 
\texttt{Always give your single{-}word clue and your comma{-}separated list of related target words in the following format and make your answers as short as possible, never include any other text than is required in this form:} \\
\\ 
\texttt{CLUE: <WORD>} \\
\texttt{TARGETS: <TARGETS>} \\
\\ 
\texttt{Your team words are: chilliness, defacement, liquidness, glasshouse, eggroll, ungodliness, bandmaster, unbefitting, ninepin.} \\
\texttt{Your opponent's team words are: secession, gentlemanliness, candleholder, shedder, inexistence, plantain, hydroplane, hotelkeeper.} \\
\texttt{Distractor words are: slouchy, musicologist, cattiness, duplicator, duckweed, bagpiper, toxicologist.} \\
\texttt{The kill words are: geophysics.} \\
            }
        }
    }
    & \\ \\

    \theutterance \stepcounter{utterance}  
    & \multicolumn{4}{p{0.6\linewidth}}{
        \cellcolor[rgb]{0.9,0.9,0.9}{
            \makecell[{{p{\linewidth}}}]{
                \texttt{\tiny{[P1$\rangle$GM]}}
                \texttt{<reasoning>} \\
\texttt{Let's tackle this clue{-}giving problem. The user's team has nine words: chilliness, defacement, liquidness, glasshouse, eggroll, ungodliness, bandmaster, unbefitting, ninepin. I need to find a single{-}word clue that relates to as many of these as possible without being one of the words themselves or containing parts of them.} \\
\\ 
\texttt{First, I'll look for common themes or relationships among the words. Let's list them out again:} \\
\\ 
\texttt{Chilliness – relates to cold, temperature.} \\
\texttt{Defacement – damaging or spoiling something.} \\
\texttt{Liquidness – state of being liquid.} \\
\texttt{Glasshouse – a greenhouse made of glass.} \\
\texttt{Eggroll – a food item, rolled egg dish.} \\
\texttt{Ungodliness – lack of morality, sinfulness.} \\
\texttt{Bandmaster – leader of a band.} \\
\texttt{Unbefitting – not suitable or appropriate.} \\
\texttt{Ninepin – a type of bowling pin.} \\
\\ 
\texttt{Hmm, not immediately obvious. Let's see if there are any overlaps. Maybe looking for parts of words or synonyms. For example, "glasshouse" and "eggroll" – maybe "house" and "roll"? Not sure. "Liquidness" and "glasshouse" – maybe "glass" is part of "glasshouse" but "liquidness" is about liquid. Not directly related.} \\
\\ 
\texttt{Wait, "eggroll" and "ninepin" – both have "roll" in them? No, ninepin is a bowling term. Eggroll is a food. Maybe "roll" isn't the connection. Let's think of other angles.} \\
\\ 
\texttt{Looking at "defacement" and "ungodliness" – both involve negative actions or states. "Defacement" is damaging, "ungodliness" is immoral. Maybe "sin" or "damage"? But "sin" might relate to "ungodliness" but not the others. "Damage" relates to "defacement" but not the rest.} \\
\\ 
\texttt{"Bandmaster" – leader of a band. Maybe "leader" or "conduct"? Not sure. "Unbefitting" – not suitable. Maybe "improper" or "unsuitable"? But how does that connect to others?} \\
\\ 
\texttt{"Chilliness" – cold. "Liquidness" – state of being liquid. Maybe "physical properties"? But that's too vague. "Glasshouse" – structure. "Eggroll" – food. "Ninepin" – object. Not seeing a clear category.} \\
\\ 
\texttt{Wait} \\
            }
        }
    }
    & & \\ \\

    \theutterance \stepcounter{utterance}  
    & & & \multicolumn{2}{p{0.3\linewidth}}{
        \cellcolor[rgb]{0.9,0.9,0.9}{
            \makecell[{{p{\linewidth}}}]{
                \texttt{\tiny{[GM$|$GM]}}
                {'player': 'Cluegiver', 'type': 'rambling error', 'utterance': '<reasoning>\nLet\'s tackle this clue-giving problem. The user\'s team has nine words: chilliness, defacement, liquidness, glasshouse, eggroll, ungodliness, bandmaster, unbefitting, ninepin. I need to find a single-word clue that relates to as many of these as possible without being one of the words themselves or containing parts of them.\n\nFirst, I\'ll look for common themes or relationships among the words. Let\'s list them out again:\n\nChilliness – relates to cold, temperature.\nDefacement – damaging or spoiling something.\nLiquidness – state of being liquid.\nGlasshouse – a greenhouse made of glass.\nEggroll – a food item, rolled egg dish.\nUngodliness – lack of morality, sinfulness.\nBandmaster – leader of a band.\nUnbefitting – not suitable or appropriate.\nNinepin – a type of bowling pin.\n\nHmm, not immediately obvious. Let\'s see if there are any overlaps. Maybe looking for parts of words or synonyms. For example, "glasshouse" and "eggroll" – maybe "house" and "roll"? Not sure. "Liquidness" and "glasshouse" – maybe "glass" is part of "glasshouse" but "liquidness" is about liquid. Not directly related.\n\nWait, "eggroll" and "ninepin" – both have "roll" in them? No, ninepin is a bowling term. Eggroll is a food. Maybe "roll" isn\'t the connection. Let\'s think of other angles.\n\nLooking at "defacement" and "ungodliness" – both involve negative actions or states. "Defacement" is damaging, "ungodliness" is immoral. Maybe "sin" or "damage"? But "sin" might relate to "ungodliness" but not the others. "Damage" relates to "defacement" but not the rest.\n\n"Bandmaster" – leader of a band. Maybe "leader" or "conduct"? Not sure. "Unbefitting" – not suitable. Maybe "improper" or "unsuitable"? But how does that connect to others?\n\n"Chilliness" – cold. "Liquidness" – state of being liquid. Maybe "physical properties"? But that\'s too vague. "Glasshouse" – structure. "Eggroll" – food. "Ninepin" – object. Not seeing a clear category.\n\nWait', 'message': 'Your answer contained more than two lines, please only give one clue and your targets on two separate lines.'}
            }
        }
    }
    & & \\ \\

\end{supertabular}
}

\end{document}
